\documentclass[11pt,a4paper]{article}
\usepackage[utf8]{inputenc}
\usepackage[T1]{fontenc}
\usepackage[catalan]{babel}
\usepackage{lmodern}
\usepackage[sfdefault,lining]{FiraSans}
\usepackage[a4paper,top=2.3cm,bottom=2.2cm,left=1.8cm,right=1.8cm,headheight=1.3cm,headsep=0.4cm,footskip=1cm]{geometry}
\usepackage{amsmath,amssymb}
\usepackage{icomma}
\usepackage{graphicx}
\usepackage[table]{xcolor}
\usepackage{tikz}
\usepackage{fontawesome5}
\usepackage{ccicons}
\usepackage{fancyhdr}
\usepackage{lastpage}
\usepackage{tasks}
\usepackage{needspace}
\usepackage[most]{tcolorbox}
\usepackage{array}
\usepackage[hidelinks,pdftitle={Fitxa d'exercicis de la unitat 4: Equacions de 1r i 2n grau. 2n d'ESO},pdfsubject={Matemàtiques 2n ESO},pdfkeywords={equacions, identitats, equacions de primer grau, parèntesis, denominadors, problemes, equacions de segon grau, incompletes, fórmula, discriminant}]{hyperref}

\definecolor{navy}{HTML}{1B2748}
\definecolor{accent}{HTML}{E8603A}
\definecolor{teal}{HTML}{1F9E90}
\definecolor{soft}{HTML}{EEF1F8}
\definecolor{muted}{HTML}{5D6685}

\setlength{\parindent}{0pt}
\setlength{\parskip}{3pt}
\settasks{label=\alph*), label-format=\bfseries\color{navy}, label-width=1.3em, item-indent=1.8em, column-sep=0.6em, after-item-skip=5pt, before-skip=3pt, after-skip=4pt}

% ---------- capçalera i peu ----------
\newcommand{\llicencia}{\href{https://creativecommons.org/licenses/by-nc-sa/4.0/deed.ca}{\ccbyncsa}\ \ Llicència Creative Commons Reconeixement-NoComercial-CompartirIgual 4.0 Internacional (CC BY-NC-SA 4.0)}
\pagestyle{fancy}
\fancyhf{}
\fancyhead[L]{\raisebox{-0.35\height}{\includegraphics[height=0.95cm]{logo-gva.png}}}
\fancyhead[C]{\small\color{muted}Fitxa d'exercicis · Unitat 4: Equacions de 1r i 2n grau · 2n d'ESO}
\fancyhead[R]{\raisebox{-0.35\height}{\includegraphics[height=0.95cm]{logo-vermellar.png}}}
\fancyfoot[L]{\scriptsize\color{muted}\llicencia}
\fancyfoot[R]{\small\color{muted}\thepage}
\renewcommand{\headrulewidth}{0.4pt}
\renewcommand{\headrule}{\color{muted!40}\hrule height \headrulewidth}
\fancypagestyle{primera}{\fancyhf{}\renewcommand{\headrulewidth}{0pt}
  \fancyfoot[L]{\scriptsize\color{muted}\llicencia}
  \fancyfoot[R]{\small\color{muted}\thepage}}

% ---------- nivells i exercicis ----------
\newcommand{\estrelles}[1]{{\color{accent}\ifcase#1\or\faStar\faStar[regular]\faStar[regular]\or\faStar\faStar\faStar[regular]\or\faStar\faStar\faStar\fi}}
\newcounter{exer}
\newcommand{\ex}[1]{\par\needspace{4\baselineskip}\medskip\stepcounter{exer}\noindent{\color{navy}\bfseries\theexer.}\ \ {\small\estrelles{#1}}\ \ }

\newcommand{\seccio}[2]{\par\needspace{10\baselineskip}\bigskip
  \begin{tcolorbox}[colback=navy,colframe=navy,coltext=white,arc=2mm,boxrule=0pt,left=3mm,right=3mm,top=1.2mm,bottom=1.2mm]
  {\large\bfseries #1\quad #2}
  \end{tcolorbox}\vspace{-2pt}}

\newcolumntype{C}[1]{>{\centering\arraybackslash}p{#1}}
% quadre de recordatori (només a la fitxa d'exercicis)
\newcommand{\recorda}[1]{\par\smallskip\begin{tcolorbox}[colback=accent!7,colframe=accent!55,arc=2mm,boxrule=0.5pt,left=2.5mm,right=2.5mm,top=1mm,bottom=1mm]\small{\color{accent}\faLightbulb[regular]}\ \textbf{Recorda:}\ #1\end{tcolorbox}}
\newcommand{\res}[1]{#1}\newcommand{\rest}[1]{#1}

\def\fmneed{12}
% =====================================================================
% Capçalera FentMates de les fitxes: un color per a cada curs.
% Els colors són els de la pàgina inicial (LEVELS a assets/js/materials.js);
% el logo ve de _eines/logo/pdf/ (python3 _eines/logo/genera.py).
%
% Ús, al final del preàmbul de la unitat:
%   % =====================================================================
% Capçalera FentMates de les fitxes: un color per a cada curs.
% Els colors són els de la pàgina inicial (LEVELS a assets/js/materials.js);
% el logo ve de _eines/logo/pdf/ (python3 _eines/logo/genera.py).
%
% Ús, al final del preàmbul de la unitat:
%   % =====================================================================
% Capçalera FentMates de les fitxes: un color per a cada curs.
% Els colors són els de la pàgina inicial (LEVELS a assets/js/materials.js);
% el logo ve de _eines/logo/pdf/ (python3 _eines/logo/genera.py).
%
% Ús, al final del preàmbul de la unitat:
%   \input{../comu/fentmates.tex}
%   \fmcurs{1eso}
% i, al començament del document:
%   \capcalera{Matemàtiques}{1r ESO}{Fitxa d'exercicis · Unitat 2}{Divisibilitat}
% Opcional abans de \input: \def\fmneed{14} (espai mínim abans de cada secció).
% =====================================================================
\makeatletter
\def\fm@logodir{../../logo/pdf}
\providecommand{\fmneed}{10}
\def\fm@cursnom{}

% clau del curs, color i color del text damunt del color
\newcommand{\fm@curs}[3]{\@namedef{fm@color@#1}{#2}\@namedef{fm@text@#1}{#3}}
\fm@curs{1eso}{E6B84A}{navy}
\fm@curs{2eso}{3B82F6}{white}
\fm@curs{3eso}{FF7A50}{white}
\fm@curs{4eso-a}{A78BFA}{white}
\fm@curs{4eso-b}{6D28D9}{white}
\fm@curs{1batx-mi}{9B1C1C}{white}
\fm@curs{1batx-ccss}{EF6B6B}{white}
\fm@curs{2batx-mii}{1E6B47}{white}
\fm@curs{2batx-ccss}{2BB3A3}{white}

\newcommand{\fmcurs}[1]{%
  \@ifundefined{fm@color@#1}{\PackageError{fentmates}{Curs desconegut: #1}{}}{}%
  \edef\fm@tmp{\noexpand\definecolor{curs}{HTML}{\@nameuse{fm@color@#1}}%
               \noexpand\colorlet{curstext}{\@nameuse{fm@text@#1}}}\fm@tmp
  \colorlet{cursfosc}{curs!72!black}%
  \def\fm@logo{\fm@logodir/vertical-#1.pdf}%
  \def\fm@marca{\fm@logodir/marca-#1.pdf}}

% ---------- capçalera de la primera pàgina ----------
% #1 matèria (petit), #2 curs, #3 tipus de document i unitat, #4 títol
\newcommand{\capcalera}[4]{%
  \gdef\fm@cursnom{#2}%
  \par\noindent
  \begin{tcolorbox}[colback=curs,colframe=curs,boxrule=0pt,arc=2.2mm,width=3.5cm,height=2.15cm,
      valign=center,left=3mm,right=2mm,top=1mm,bottom=1mm,nobeforeafter,box align=center]
    \color{curstext}{\footnotesize #1}\par\vspace{1pt}{\fontsize{19}{21}\selectfont\bfseries #2\par}%
  \end{tcolorbox}\hspace{5mm}%
  \begin{minipage}[c]{\dimexpr\linewidth-3.5cm-5mm-2.9cm\relax}
    {\color{cursfosc}\small\bfseries\MakeUppercase{#3}}\par\vspace{2pt}
    {\color{navy}\fontsize{20}{23}\selectfont\bfseries #4\par}
  \end{minipage}\hfill
  \begin{minipage}[c]{2.6cm}\raggedleft\includegraphics[width=2.6cm]{\fm@logo}\end{minipage}\par
  \vspace{2pt}}

% ---------- barres de secció del color del curs ----------
\renewcommand{\seccio}[2]{\par\needspace{\fmneed\baselineskip}\bigskip
  \begin{tcolorbox}[colback=curs,colframe=curs,coltext=curstext,arc=2mm,boxrule=0pt,left=3mm,right=3mm,top=1.2mm,bottom=1.2mm]
  {\large\bfseries #1\quad #2}
  \end{tcolorbox}\vspace{-2pt}}

% ---------- estrelles de dificultat del color del curs (a Batxillerat no se'n posen) ----------
\@ifundefined{estrelles}{}{%
  \renewcommand{\estrelles}[1]{{\color{curs}\ifcase#1\or\faStar\faStar[regular]\faStar[regular]\or\faStar\faStar\faStar[regular]\or\faStar\faStar\faStar\fi}}}

% ---------- capçalera i peu de les pàgines ----------
% primera pàgina: logos institucionals al peu
\fancypagestyle{primera}{\fancyhf{}\renewcommand{\headrulewidth}{0pt}
  \fancyfoot[L]{\raisebox{-0.32\height}{\includegraphics[height=0.78cm]{logo-gva.png}}\hspace{5mm}%
    \raisebox{-0.32\height}{\includegraphics[height=0.78cm]{logo-vermellar.png}}}
  \fancyfoot[R]{\scriptsize\color{muted}\href{https://creativecommons.org/licenses/by-nc-sa/4.0/deed.ca}{\ccbyncsa}\ \ CC BY-NC-SA 4.0 · fentmates.org\qquad\small\thepage}}
% la resta: el símbol a l'esquerra i el curs a la dreta, amb el títol al centre com abans
\fancyhead[L]{\raisebox{-0.3\height}{\includegraphics[height=0.62cm]{\fm@marca}}}
\fancyhead[R]{\small\bfseries\color{cursfosc}\fm@cursnom}
\renewcommand{\headrule}{\color{curs!45}\hrule height \headrulewidth}
\makeatother

%   \fmcurs{1eso}
% i, al començament del document:
%   \capcalera{Matemàtiques}{1r ESO}{Fitxa d'exercicis · Unitat 2}{Divisibilitat}
% Opcional abans de \input: \def\fmneed{14} (espai mínim abans de cada secció).
% =====================================================================
\makeatletter
\def\fm@logodir{../../logo/pdf}
\providecommand{\fmneed}{10}
\def\fm@cursnom{}

% clau del curs, color i color del text damunt del color
\newcommand{\fm@curs}[3]{\@namedef{fm@color@#1}{#2}\@namedef{fm@text@#1}{#3}}
\fm@curs{1eso}{E6B84A}{navy}
\fm@curs{2eso}{3B82F6}{white}
\fm@curs{3eso}{FF7A50}{white}
\fm@curs{4eso-a}{A78BFA}{white}
\fm@curs{4eso-b}{6D28D9}{white}
\fm@curs{1batx-mi}{9B1C1C}{white}
\fm@curs{1batx-ccss}{EF6B6B}{white}
\fm@curs{2batx-mii}{1E6B47}{white}
\fm@curs{2batx-ccss}{2BB3A3}{white}

\newcommand{\fmcurs}[1]{%
  \@ifundefined{fm@color@#1}{\PackageError{fentmates}{Curs desconegut: #1}{}}{}%
  \edef\fm@tmp{\noexpand\definecolor{curs}{HTML}{\@nameuse{fm@color@#1}}%
               \noexpand\colorlet{curstext}{\@nameuse{fm@text@#1}}}\fm@tmp
  \colorlet{cursfosc}{curs!72!black}%
  \def\fm@logo{\fm@logodir/vertical-#1.pdf}%
  \def\fm@marca{\fm@logodir/marca-#1.pdf}}

% ---------- capçalera de la primera pàgina ----------
% #1 matèria (petit), #2 curs, #3 tipus de document i unitat, #4 títol
\newcommand{\capcalera}[4]{%
  \gdef\fm@cursnom{#2}%
  \par\noindent
  \begin{tcolorbox}[colback=curs,colframe=curs,boxrule=0pt,arc=2.2mm,width=3.5cm,height=2.15cm,
      valign=center,left=3mm,right=2mm,top=1mm,bottom=1mm,nobeforeafter,box align=center]
    \color{curstext}{\footnotesize #1}\par\vspace{1pt}{\fontsize{19}{21}\selectfont\bfseries #2\par}%
  \end{tcolorbox}\hspace{5mm}%
  \begin{minipage}[c]{\dimexpr\linewidth-3.5cm-5mm-2.9cm\relax}
    {\color{cursfosc}\small\bfseries\MakeUppercase{#3}}\par\vspace{2pt}
    {\color{navy}\fontsize{20}{23}\selectfont\bfseries #4\par}
  \end{minipage}\hfill
  \begin{minipage}[c]{2.6cm}\raggedleft\includegraphics[width=2.6cm]{\fm@logo}\end{minipage}\par
  \vspace{2pt}}

% ---------- barres de secció del color del curs ----------
\renewcommand{\seccio}[2]{\par\needspace{\fmneed\baselineskip}\bigskip
  \begin{tcolorbox}[colback=curs,colframe=curs,coltext=curstext,arc=2mm,boxrule=0pt,left=3mm,right=3mm,top=1.2mm,bottom=1.2mm]
  {\large\bfseries #1\quad #2}
  \end{tcolorbox}\vspace{-2pt}}

% ---------- estrelles de dificultat del color del curs (a Batxillerat no se'n posen) ----------
\@ifundefined{estrelles}{}{%
  \renewcommand{\estrelles}[1]{{\color{curs}\ifcase#1\or\faStar\faStar[regular]\faStar[regular]\or\faStar\faStar\faStar[regular]\or\faStar\faStar\faStar\fi}}}

% ---------- capçalera i peu de les pàgines ----------
% primera pàgina: logos institucionals al peu
\fancypagestyle{primera}{\fancyhf{}\renewcommand{\headrulewidth}{0pt}
  \fancyfoot[L]{\raisebox{-0.32\height}{\includegraphics[height=0.78cm]{logo-gva.png}}\hspace{5mm}%
    \raisebox{-0.32\height}{\includegraphics[height=0.78cm]{logo-vermellar.png}}}
  \fancyfoot[R]{\scriptsize\color{muted}\href{https://creativecommons.org/licenses/by-nc-sa/4.0/deed.ca}{\ccbyncsa}\ \ CC BY-NC-SA 4.0 · fentmates.org\qquad\small\thepage}}
% la resta: el símbol a l'esquerra i el curs a la dreta, amb el títol al centre com abans
\fancyhead[L]{\raisebox{-0.3\height}{\includegraphics[height=0.62cm]{\fm@marca}}}
\fancyhead[R]{\small\bfseries\color{cursfosc}\fm@cursnom}
\renewcommand{\headrule}{\color{curs!45}\hrule height \headrulewidth}
\makeatother

%   \fmcurs{1eso}
% i, al començament del document:
%   \capcalera{Matemàtiques}{1r ESO}{Fitxa d'exercicis · Unitat 2}{Divisibilitat}
% Opcional abans de \input: \def\fmneed{14} (espai mínim abans de cada secció).
% =====================================================================
\makeatletter
\def\fm@logodir{../../logo/pdf}
\providecommand{\fmneed}{10}
\def\fm@cursnom{}

% clau del curs, color i color del text damunt del color
\newcommand{\fm@curs}[3]{\@namedef{fm@color@#1}{#2}\@namedef{fm@text@#1}{#3}}
\fm@curs{1eso}{E6B84A}{navy}
\fm@curs{2eso}{3B82F6}{white}
\fm@curs{3eso}{FF7A50}{white}
\fm@curs{4eso-a}{A78BFA}{white}
\fm@curs{4eso-b}{6D28D9}{white}
\fm@curs{1batx-mi}{9B1C1C}{white}
\fm@curs{1batx-ccss}{EF6B6B}{white}
\fm@curs{2batx-mii}{1E6B47}{white}
\fm@curs{2batx-ccss}{2BB3A3}{white}

\newcommand{\fmcurs}[1]{%
  \@ifundefined{fm@color@#1}{\PackageError{fentmates}{Curs desconegut: #1}{}}{}%
  \edef\fm@tmp{\noexpand\definecolor{curs}{HTML}{\@nameuse{fm@color@#1}}%
               \noexpand\colorlet{curstext}{\@nameuse{fm@text@#1}}}\fm@tmp
  \colorlet{cursfosc}{curs!72!black}%
  \def\fm@logo{\fm@logodir/vertical-#1.pdf}%
  \def\fm@marca{\fm@logodir/marca-#1.pdf}}

% ---------- capçalera de la primera pàgina ----------
% #1 matèria (petit), #2 curs, #3 tipus de document i unitat, #4 títol
\newcommand{\capcalera}[4]{%
  \gdef\fm@cursnom{#2}%
  \par\noindent
  \begin{tcolorbox}[colback=curs,colframe=curs,boxrule=0pt,arc=2.2mm,width=3.5cm,height=2.15cm,
      valign=center,left=3mm,right=2mm,top=1mm,bottom=1mm,nobeforeafter,box align=center]
    \color{curstext}{\footnotesize #1}\par\vspace{1pt}{\fontsize{19}{21}\selectfont\bfseries #2\par}%
  \end{tcolorbox}\hspace{5mm}%
  \begin{minipage}[c]{\dimexpr\linewidth-3.5cm-5mm-2.9cm\relax}
    {\color{cursfosc}\small\bfseries\MakeUppercase{#3}}\par\vspace{2pt}
    {\color{navy}\fontsize{20}{23}\selectfont\bfseries #4\par}
  \end{minipage}\hfill
  \begin{minipage}[c]{2.6cm}\raggedleft\includegraphics[width=2.6cm]{\fm@logo}\end{minipage}\par
  \vspace{2pt}}

% ---------- barres de secció del color del curs ----------
\renewcommand{\seccio}[2]{\par\needspace{\fmneed\baselineskip}\bigskip
  \begin{tcolorbox}[colback=curs,colframe=curs,coltext=curstext,arc=2mm,boxrule=0pt,left=3mm,right=3mm,top=1.2mm,bottom=1.2mm]
  {\large\bfseries #1\quad #2}
  \end{tcolorbox}\vspace{-2pt}}

% ---------- estrelles de dificultat del color del curs (a Batxillerat no se'n posen) ----------
\@ifundefined{estrelles}{}{%
  \renewcommand{\estrelles}[1]{{\color{curs}\ifcase#1\or\faStar\faStar[regular]\faStar[regular]\or\faStar\faStar\faStar[regular]\or\faStar\faStar\faStar\fi}}}

% ---------- capçalera i peu de les pàgines ----------
% primera pàgina: logos institucionals al peu
\fancypagestyle{primera}{\fancyhf{}\renewcommand{\headrulewidth}{0pt}
  \fancyfoot[L]{\raisebox{-0.32\height}{\includegraphics[height=0.78cm]{logo-gva.png}}\hspace{5mm}%
    \raisebox{-0.32\height}{\includegraphics[height=0.78cm]{logo-vermellar.png}}}
  \fancyfoot[R]{\scriptsize\color{muted}\href{https://creativecommons.org/licenses/by-nc-sa/4.0/deed.ca}{\ccbyncsa}\ \ CC BY-NC-SA 4.0 · fentmates.org\qquad\small\thepage}}
% la resta: el símbol a l'esquerra i el curs a la dreta, amb el títol al centre com abans
\fancyhead[L]{\raisebox{-0.3\height}{\includegraphics[height=0.62cm]{\fm@marca}}}
\fancyhead[R]{\small\bfseries\color{cursfosc}\fm@cursnom}
\renewcommand{\headrule}{\color{curs!45}\hrule height \headrulewidth}
\makeatother

\fmcurs{2eso}

\begin{document}
\thispagestyle{primera}

\capcalera{Matemàtiques}{2n ESO}{Fitxa d'exercicis · Unitat 4}{Equacions de 1r i 2n grau}
\vspace{10pt}
Nom i cognoms: \rule{8.4cm}{0.4pt}\hfill Grup: \rule{1.6cm}{0.4pt}\hfill Data: \rule{2.2cm}{0.4pt}
\vspace{8pt}

\begin{tcolorbox}[colback=soft,colframe=soft,arc=2mm,left=3mm,right=3mm,top=1.5mm,bottom=1.5mm]
\small Exercicis graduats per nivells de dificultat:\quad \estrelles{1}\ bàsic \qquad \estrelles{2}\ mitjà \qquad \estrelles{3}\ avançat.\\[2pt] Resol-los al quadern mostrant tots els passos i comprova les solucions sempre que pugues. Fes els càlculs sense calculadora. Les seccions segueixen l'ordre dels apunts de la unitat (4.1--4.9).
\end{tcolorbox}

% =====================================================
\seccio{1}{Equacions: elements i solucions}

\ex{1} Indica quines igualtats són identitats i quines són equacions:\begin{tasks}(2)
  \task $3(x+2)=3x+6$
  \task $3x+2=8$
  \task $(x+1)^{2}=x^{2}+2x+1$
  \task $x^{2}-1=0$
  \task $5x-2x=3x$
  \task $5x-2=3x$
  \task $(x-3)(x+3)=x^{2}-9$
  \task $2(x-1)=2x-1$
\end{tasks}\recorda{una \textbf{identitat} es compleix per a qualsevol valor de $x$; una \textbf{equació}, només per a alguns valors (o per a cap).}

\ex{1} Comprova si el valor indicat és solució de cada equació:\begin{tasks}(2)
  \task $3x-5=x+1$, \ $x=3$
  \task $x^{2}-4x=5$, \ $x=-1$
  \task $2(x-1)=x+4$, \ $x=5$
  \task $\dfrac{x}{2}+1=x-3$, \ $x=8$
  \task $x^{2}+x=6$, \ $x=-3$
  \task $7-x=2x+1$, \ $x=3$
\end{tasks}

\ex{1} En l'equació $4x^2+x-3=2x+5$, indica el primer i el segon membre, els termes, la incògnita i el grau.

\ex{2} Indica el grau de cada equació un cop reduïda:\begin{tasks}(2)
  \task $x(x+4)=x^{2}-8$
  \task $(x-2)^{2}=x^{2}+4$
  \task $x^{2}+3x=2x^{2}-1$
  \task $(x+1)(x-1)=3x$
\end{tasks}\recorda{abans de dir el grau, opera i redueix: els termes en $x^2$ es poden cancel·lar.}

\ex{2} Escriu una equació de primer grau que tinga la solució $x=-2$ i una equació de segon grau que tinga les solucions $x=3$ i $x=-3$.

\ex{3} Troba el valor de $k$ en cada cas:\begin{tasks}(1)
  \task $x=3$ és solució de $kx-4=2x+k$.
  \task $x=-1$ és solució de $x^2+kx=k+5$.
\end{tasks}


% =====================================================
\seccio{2}{Equacions de primer grau}

\ex{1} Resol mentalment o en un sol pas:\begin{tasks}(4)
  \task $x+7=12$
  \task $x-5=-3$
  \task $3x=-21$
  \task $-4x=20$
  \task $\dfrac{x}{5}=3$
  \task $-x=6$
  \task $2x=7$
  \task $\dfrac{x}{4}=-2$
\end{tasks}\recorda{el que suma passa restant i el que resta passa sumant; el que multiplica tot un membre passa dividint (amb el seu signe) i el que divideix passa multiplicant.}

\ex{1} Resol:\begin{tasks}(3)
  \task $5x-4=11$
  \task $-2x+9=1$
  \task $3x+8=2$
  \task $7-x=10$
  \task $4x-3=-15$
  \task $6+5x=-4$
\end{tasks}

\ex{1} Resol:\begin{tasks}(3)
  \task $4x+3=2x+11$
  \task $7-3x=x-13$
  \task $6x-5=9x+7$
  \task $2x-1=5x-1$
  \task $8x+2=3x-13$
  \task $10-x=4x$
\end{tasks}

\ex{2} Redueix cada membre i resol:\begin{tasks}(2)
  \task $5x-3+2x=4x+9$
  \task $8-3x+x=2x-4-x$
  \task $2x+7-5x=x-9$
  \task $x-4+3x=6x+2-x$
\end{tasks}

\ex{2} Resol (les solucions són fraccions: simplifica-les):\begin{tasks}(4)
  \task $6x+1=2x+7$
  \task $5x-2=2x+3$
  \task $4-7x=x$
  \task $9x+4=3x-6$
\end{tasks}

\ex{3} Resol i comprova la solució:\begin{tasks}(2)
  \task $3x-8+5x-1=4x+7-2x+2$
  \task $12-5x+3x=2x-9+x+1$
  \task $7x-2-x+5=4x+11-3x$
  \task $-x+6-4x=2-3x+x-10$
\end{tasks}

\ex{3} L'Aitana ha resolt l'equació $5x-3=2x+9$ així:\quad $5x+2x=9-3\ \Rightarrow\ 7x=6\ \Rightarrow\ x=\tfrac67$. \ Quins errors ha comés? Resol-la correctament i comprova-la.


% =====================================================
\seccio{3}{Equacions amb parèntesis}

\ex{1} Resol traient primer els parèntesis:\begin{tasks}(2)
  \task $3(x-2)=12$
  \task $5(x+1)=2x-7$
  \task $2(x-4)=x+1$
  \task $4(2x-1)=3x+6$
\end{tasks}

\ex{1} Resol:\begin{tasks}(4)
  \task $4-(x+3)=2x-5$
  \task $x-(5-x)=3$
  \task $7-(2x-1)=x+2$
  \task $-(x-6)=2x$
\end{tasks}\recorda{un signe $-$ davant d'un parèntesi canvia el signe de \textbf{tots} els termes de dins: $-(2x-1)=-2x+1$.}

\ex{2} Resol:\begin{tasks}(2)
  \task $3(2x-1)-2(x+4)=5x-13$
  \task $5-2(3-x)=4(x+1)-11$
  \task $4(x-3)-(x-5)=2(x+1)$
  \task $2(x+3)-3(x-1)=4x-6$
\end{tasks}

\ex{2} Resol i digues quines equacions no tenen solució i quines en tenen infinites:\begin{tasks}(2)
  \task $3(x+2)-x=2(x+3)$
  \task $5x-2(x-1)=3x+4$
  \task $6x-2(3x+1)=-2$
  \task $7x-3(2x-1)=2x+1$
\end{tasks}

\ex{2} Resol fent servir les identitats notables (veuràs que els termes en $x^2$ es cancel·len):\begin{tasks}(3)
  \task $(x+1)^{2}-(x-1)^{2}=12$
  \task $x(x+3)-x^{2}=2x+5$
  \task $(x+2)(x-2)=x(x+1)-7$
\end{tasks}

\ex{3} Resol (comença pels parèntesis de dins):\begin{tasks}(2)
  \task $x-[2x-(3-x)]=7$
  \task $2[3-2(x+1)]=5-3(2x-1)-x$
\end{tasks}

\ex{3} Troba el valor de $a$ perquè l'equació $2(x+a)=x+9$ tinga la solució $x=3$.


% =====================================================
\seccio{4}{Equacions amb denominadors}

\ex{1} Resol:\begin{tasks}(4)
  \task $\dfrac{x}{3}=4$
  \task $\dfrac{x}{2}+\dfrac{x}{4}=6$
  \task $x-\dfrac{x}{3}=10$
  \task $\dfrac{2x}{5}-1=3$
\end{tasks}

\ex{1} Resol:\begin{tasks}(4)
  \task $\dfrac{x+2}{3}=4$
  \task $\dfrac{x-1}{2}=\dfrac{x+3}{4}$
  \task $\dfrac{3x-1}{4}=2$
  \task $\dfrac{x}{6}+1=\dfrac{x}{4}$
\end{tasks}\recorda{multiplica \textbf{tots} els termes pel m.c.m. dels denominadors (també els que no tenen denominador) i posa els numeradors entre parèntesis.}

\ex{2} Resol:\begin{tasks}(2)
  \task $\dfrac{x+1}{2}-\dfrac{x-1}{3}=2$
  \task $\dfrac{x-2}{4}-\dfrac{x-3}{6}=\dfrac{1}{2}$
  \task $\dfrac{2x+1}{3}-\dfrac{x-2}{2}=3$
  \task $x-\dfrac{x+2}{5}=\dfrac{x+6}{2}$
\end{tasks}

\ex{2} Resol:\begin{tasks}(2)
  \task $\dfrac{x-1}{3}+\dfrac{x+2}{6}=x-2$
  \task $3-\dfrac{x-1}{4}=\dfrac{x+5}{2}$
\end{tasks}

\ex{3} Resol i interpreta el resultat:\begin{tasks}(1)
  \task $\dfrac{3(x-2)}{4}-\dfrac{2(x+1)}{3}=\dfrac{x-10}{6}$
  \task $\dfrac{x}{2}-\dfrac{x-1}{3}-\dfrac{x+1}{4}=\dfrac{1-x}{12}$
  \task $\dfrac{x+1}{2}-\dfrac{x+3}{4}=\dfrac{x}{4}-1$
\end{tasks}

\ex{3} Un alumne ha resolt l'equació $\dfrac{x-2}{3}-\dfrac{x+1}{2}=1$ així:\quad $2x-4-3x+3=6\ \Rightarrow\ -x=7\ \Rightarrow\ x=-7$. \ On s'ha equivocat? Resol-la correctament.


% =====================================================
\seccio{5}{Problemes amb equacions de primer grau}

\ex{1} La suma de dos nombres consecutius és 75. Quins són?\recorda{tria la incògnita i escriu què representa, planteja l'equació, resol-la i comprova que la solució té sentit en el problema.}

\ex{1} Si a un nombre li sumem el seu triple, obtenim 56. Quin és el nombre?

\ex{1} L'Anna té 12 anys i son pare, 40. D'ací a quants anys l'edat del pare serà el doble que la de l'Anna?

\ex{2} La llargària d'un rectangle és 3 m més que el doble de l'amplària i el perímetre fa 54 m. Quines dimensions té?

\ex{2} En un corral hi ha gallines i conills. En total hi ha 40 caps i 110 potes. Quants animals hi ha de cada tipus?

\ex{2} En una guardiola hi ha 30 monedes, unes d'1 € i altres de 2 €, que sumen 46 €. Quantes monedes hi ha de cada tipus?

\ex{2} L'edat d'una mare és el triple de la del seu fill, i d'ací a 12 anys només serà el doble. Quants anys tenen ara?

\ex{2} Després d'un descompte del 15 \%, un abric costa 68 €. Quant costava abans del descompte?

\ex{3} Una persona gasta la tercera part dels seus diners en roba i la quarta part en llibres, i encara li queden 25 €. Quants diners tenia?

\ex{3} Repartim 150 € entre tres germans: el segon rep el doble que el primer, i el tercer, 30 € més que el segon. Quant rep cadascun?

\ex{3} Volem 20 kg d'una mescla de cafè a 9,50 €/kg barrejant cafè de 12 €/kg amb cafè de 8 €/kg. Quants quilos de cada tipus hem de posar?

\ex{3} Dos ciclistes ixen alhora de dues ciutats separades 90 km i van l'un cap a l'altre a 20 km/h i a 25 km/h. Quant de temps tarden a trobar-se? A quina distància de la primera ciutat?


% =====================================================
\seccio{6}{Equacions de segon grau incompletes}

\ex{1} Escriu cada equació en la forma general $ax^2+bx+c=0$, indica els coeficients $a$, $b$ i $c$ i digues si és completa o incompleta:\begin{tasks}(2)
  \task $5x-x^{2}=6$
  \task $3x^{2}=12$
  \task $x(x-4)=0$
  \task $(x+1)^{2}=2x+5$
\end{tasks}

\ex{1} Resol:\begin{tasks}(3)
  \task $x^{2}-36=0$
  \task $2x^{2}-50=0$
  \task $x^{2}+9=0$
  \task $4x^{2}-9=0$
  \task $3x^{2}=27$
  \task $x^{2}-7=0$
\end{tasks}\recorda{en $ax^2+c=0$ aïllem $x^2$ i fem l'arrel: hi ha \textbf{dues} solucions, $x=\pm\sqrt{\ }$, o cap si el nombre és negatiu.}

\ex{1} Resol:\begin{tasks}(4)
  \task $x^{2}-5x=0$
  \task $3x^{2}+12x=0$
  \task $2x^{2}=7x$
  \task $-x^{2}+6x=0$
\end{tasks}\recorda{en $ax^2+bx=0$ traiem factor comú $x$: una solució sempre és $x=0$. No dividisques entre $x$!}

\ex{2} Redueix i resol:\begin{tasks}(2)
  \task $(x+3)(x-3)=7$
  \task $x(x-2)=4-2x$
  \task $(2x-1)^{2}=1$
  \task $(x+2)^{2}-4=3x$
\end{tasks}

\ex{2} Resol:\begin{tasks}(2)
  \task $\dfrac{x^{2}}{3}=12$
  \task $\dfrac{x^{2}}{4}-\dfrac{x}{2}=0$
  \task $\dfrac{x^{2}+1}{5}=2$
  \task $\dfrac{x^{2}-2}{2}-\dfrac{x^{2}+1}{3}=1$
\end{tasks}

\ex{3} La Marta resol $5x^2=3x$ dividint els dos membres entre $x$ i obté $5x=3$, és a dir, $x=\tfrac35$. Què ha fet malament? Resol-la correctament.

\ex{3} Escriu una equació de segon grau incompleta que tinga per solucions:\begin{tasks}(3)
  \task $4$ i $-4$
  \task $0$ i $3$
  \task $0$ i $-\tfrac12$
\end{tasks}


% =====================================================
\seccio{7}{Equacions de segon grau completes}

\ex{1} Resol amb la fórmula:\begin{tasks}(4)
  \task $x^{2}-6x+8=0$
  \task $x^{2}+2x-8=0$
  \task $x^{2}-4x-5=0$
  \task $x^{2}+7x+12=0$
\end{tasks}\recorda{$x=\dfrac{-b\pm\sqrt{b^2-4ac}}{2a}$. \ Compte amb els signes: si $b=-6$, aleshores $-b=6$ i $b^2=36$.}

\ex{1} Resol:\begin{tasks}(4)
  \task $2x^{2}-5x+2=0$
  \task $3x^{2}+x-2=0$
  \task $6x^{2}-x-1=0$
  \task $4x^{2}+4x-3=0$
\end{tasks}

\ex{2} Resol (si $a$ és negatiu, multiplica per $-1$; si tots els coeficients tenen un divisor comú, simplifica):\begin{tasks}(4)
  \task $-x^{2}+x+6=0$
  \task $2x^{2}+4x-6=0$
  \task $x^{2}-x-20=0$
  \task $3x^{2}-3x-18=0$
\end{tasks}

\ex{2} Redueix a la forma general i resol:\begin{tasks}(2)
  \task $(x-1)(x+4)=6$
  \task $(x+2)^{2}=3x+10$
  \task $x(x+3)=2(x+3)$
  \task $3x(x-2)-(x-1)(x+1)=9$
\end{tasks}

\ex{3} Resol:\begin{tasks}(2)
  \task $\dfrac{x^{2}}{2}-\dfrac{x}{3}=\dfrac{4}{3}$
  \task $\dfrac{x^{2}-1}{3}+\dfrac{x}{2}=2$
\end{tasks}

\ex{3} Resol:\begin{tasks}(2)
  \task $\dfrac{x(x-1)}{3}=x-1$
  \task $\dfrac{x(x+1)}{2}-\dfrac{(x-1)(x+2)}{3}=1$
\end{tasks}

\ex{3} Escriu una equació de segon grau, en forma general, que tinga per solucions:\begin{tasks}(2)
  \task $2$ i $-5$
  \task $\tfrac12$ i $3$
\end{tasks}\recorda{si les solucions són $r$ i $s$, l'equació $(x-r)(x-s)=0$ les té.}


% =====================================================
\seccio{8}{El discriminant}

\ex{1} Calcula el discriminant i digues quantes solucions té cada equació, sense resoldre-la:\begin{tasks}(4)
  \task $x^{2}-10x+25=0$
  \task $2x^{2}+3x+4=0$
  \task $x^{2}-x-1=0$
  \task $-3x^{2}+2x+1=0$
\end{tasks}\recorda{$\Delta=b^2-4ac$. \ Si $\Delta>0$, dues solucions; si $\Delta=0$, una (doble); si $\Delta<0$, cap.}

\ex{1} Quantes solucions té cada equació?\begin{tasks}(3)
  \task $x^{2}+4x+4=0$
  \task $x^{2}+x-12=0$
  \task $5x^{2}-2x+1=0$
  \task $x^{2}-9=0$
  \task $4x^{2}-12x+9=0$
  \task $x^{2}+3x=0$
\end{tasks}

\ex{2} Resol (fixa't en el discriminant abans de continuar):\begin{tasks}(4)
  \task $9x^{2}-6x+1=0$
  \task $x^{2}+3x+5=0$
  \task $x^{2}-2x-2=0$
  \task $x^{2}+4x+1=0$
\end{tasks}

\ex{2} Escriu un valor de $c$ perquè l'equació $x^2+6x+c=0$ tinga: a) dues solucions; b) una sola solució; c) cap solució.

\ex{3} Respon:\begin{tasks}(1)
  \task Troba $k$ perquè l'equació $x^2-4x+k=0$ tinga una única solució. Quina és?
  \task Troba $m$ perquè $x^2+mx+16=0$ tinga una solució doble.
\end{tasks}

\ex{3} Respon:\begin{tasks}(1)
  \task Per a quins valors de $m$ l'equació $x^2-2x+m=0$ té dues solucions?
  \task Troba $k$ perquè $x=3$ siga solució de $x^2+kx-6=0$ i calcula l'altra solució.
\end{tasks}


% =====================================================
\seccio{9}{Problemes amb equacions de segon grau}

\ex{1} Si al quadrat d'un nombre li restem 5, obtenim 44. Quin pot ser el nombre?\recorda{l'equació pot tindre dues solucions: comprova quines tenen sentit en el problema (una longitud o una edat no poden ser negatives).}

\ex{1} L'àrea d'un quadrat és 144 cm². Quant fa el costat? I si l'àrea fóra 50 cm²?

\ex{1} Un rectangle fa 3 m més de llarg que d'ample i la seua àrea és 40 m². Quines dimensions té?

\ex{2} Troba dos nombres naturals consecutius el producte dels quals siga 156.

\ex{2} Si al quadrat d'un nombre li restem el seu triple, obtenim 40. Quin pot ser el nombre?

\ex{2} Els catets d'un triangle rectangle mesuren $x$ i $x+7$ cm, i la hipotenusa 13 cm. Quant mesuren els catets?

\ex{2} Si al quadrat de l'edat d'en Pau li restem el triple de la seua edat, obtenim 130. Quants anys té?

\ex{2} Si augmentem en 3 cm el costat d'un quadrat, la seua àrea augmenta 51 cm². Quant fa el costat?

\ex{3} La suma dels quadrats de dos nombres consecutius és 85. Quins poden ser?

\ex{3} Un rectangle té 34 m de perímetre i 60 m² d'àrea. Quines dimensions té?

\ex{3} La hipotenusa d'un triangle rectangle fa 10 cm i un catet fa 2 cm més que l'altre. Quant fan els catets?

\ex{3} Una piscina rectangular de 10 m × 6 m està envoltada per una vorera d'amplària constant. L'àrea total (piscina i vorera) és 96 m². Quina amplària té la vorera?

\vfill
\begin{center}\color{muted}\small\faCheckCircle[regular]\ \ Ànim! Comprova els teus resultats amb les eines interactives dels apunts de la unitat.\end{center}

\end{document}
